% Generato da plot_lqgi_revision.m; dati della chiusura reale.
\[\bar w_i=200,\quad Q_i=200 I_2,\quad R=0.1 I_2.\]
\[K_x=[K_{x,a}\ K_{x,b}].\]
{\small
\[K_{x,a}=\begin{bmatrix}
269.299 & 18.7569 & 20.3466 & 4.21621 & -167.136 & -13.737 & -1.23949 \\
18.5297 & 219.859 & 5.19353 & 7.65248 & -6.07863 & -142.077 & -0.661138 \\
\end{bmatrix}.\]
\[K_{x,b}=\begin{bmatrix}
-0.937049 & 11.652 & 22.2361 & 43.3764 & -0.976831 & -2.36859 & -13.6651 \\
1.43283 & -0.976831 & -2.23597 & -11.0676 & 14.435 & 28.7953 & 81.4582 \\
\end{bmatrix}.\]
\[K_i=\begin{bmatrix}
-44.6778 & 1.97253 \\
-1.97253 & -44.6778 \\
\end{bmatrix}.\]
}
Residuo relativo CARE: $\num{2.793777e-13}$; residuo del guadagno: $\num{1.3022188e-11}$; residuo della trasformazione di separazione: $\num{2.1806205e-13}$.
Parte reale dominante della chiusura completa: $\num{-0.821436817}$ s$^{-1}$.
Il massimo residuo LQGI a 8 s nei due esperimenti vale $\num{0.00017080376}$ rad; non coincide con il valore asintotico.\par
\begin{table}[htbp]\centering\small
\begin{tabular}{rrrrr}\toprule
$\bar w_i$ & $\|K_i\|_2$ & $T_s$ [s] & OS [\%] & picco $u$ [N m] \\
\midrule
20 & 14.142 & $\infty$ & 0 & 0.71121 \\
200 & 44.721 & 5.046 & 0 & 0.82048 \\
800 & 89.443 & 2.642 & 0 & 0.82184 \\
\bottomrule\end{tabular}\caption{LQGI: tre pesi integrali diretti a $Q$ e $R=0.1 I_2$ fissi. Gradini di 5 gradi, 8 s; $\infty$ indica mancato assestamento entro 8 s.}\label{tab:lqgi-wi}\end{table}
\begin{table}[htbp]\centering\small
\begin{tabular}{llrrrr}\toprule
Controllo & Disturbo & $q_1(\infty)$ & $q_2(\infty)$ & picco $u$ & RMS $u$ \\
\midrule
LQG & $d_1$ & 0.0718909 & -0.0280401 & 1.2925 & 0.784522 \\
LQG & $d_2$ & -0.0269505 & 0.154304 & 1.15671 & 0.571555 \\
LQGI & $d_1$ & 0 & 0 & 1.29841 & 0.447108 \\
LQGI & $d_2$ & 0 & 0 & 1.18248 & 0.394729 \\
\bottomrule\end{tabular}\caption{Disturbi separati da 0.5 N m; posizioni DC in rad, coppie in N m. Picco e RMS su entrambi i comandi negli stessi 8 s.}\label{tab:lqgi-disturbi}\end{table}
